\section{A periodic hierarchy from the physical local interactions}
\label{sec:hastings05_physical_hierarchy}

Theorem~\ref{thm:hastings05_physical_hierarchy} constructs the periodic
hierarchy from the original interactions, under the local residual
assumption used at the start of
\cite[Section~III, lines 493--507]{Hastings2006Solving}. The projection
radius is explicit and grows quadratically in the logarithm of the chain
length. The original Hamiltonian supplies the gap used to bound the
dimension of the final retained space.

\begin{lemma}[Grouping a physical cycle]
    \label{lem:hastings05_rounded_pairs}
    \lean{TNLean.MPS.AreaLaw.LocalSolution.roundedPeriodicBlock,
        TNLean.MPS.AreaLaw.LocalSolution.roundedPeriodicEdge,
        TNLean.MPS.AreaLaw.LocalSolution.roundedPeriodicDistance,
        TNLean.MPS.AreaLaw.LocalSolution.roundedPeriodicBlockCount,
        TNLean.MPS.AreaLaw.LocalSolution.periodicPhysicalMetricSpace,
        TNLean.MPS.AreaLaw.LocalSolution.card_periodicPhysical_metricRangeBall_le,
        TNLean.MPS.AreaLaw.LocalSolution.roundedPeriodicOperator,
        TNLean.MPS.AreaLaw.LocalSolution.card_roundedPeriodicEdge_fibre_le,
        TNLean.MPS.AreaLaw.LocalSolution.sum_roundedPeriodicOperator,
        TNLean.MPS.AreaLaw.LocalSolution.roundedPeriodicOperator_mem_supportedOperators,
        TNLean.MPS.AreaLaw.LocalSolution.residual_roundedPeriodicOperator_le,
        TNLean.MPS.AreaLaw.LocalSolution.centeredNorm_roundedPeriodicOperator_le}
    \leanok
    Let $M\ge1$ and $\ell\ge1$ be integers, and put
    $N=\max(1,\lfloor M/(2\ell)\rfloor)$. Divide the cycle into $N$
    consecutive blocks, including the remainder in the final block.
    Write $b(v)=\min(\lfloor v/(2\ell)\rfloor,N-1)$ and assign centre
    $i$ to edge $e(i)=\lfloor(i+\ell)/(2\ell)\rfloor$ if this is less
    than $N$, and to edge zero otherwise. For Hermitian operators $a_i$ supported
    in the cyclic ball of radius $\ell$, define
    $B_j=\sum_{e(i)=j}a_i$. Each nonclosing $B_j$ is supported on
    blocks $j-1,j$, each edge receives at most $4\ell$ terms, and
    $\sum_jB_j=\sum_i a_i$. If $\|a_i\|\le J$ and
    $\|(a_i-\langle\psi,a_i\psi\rangle)\psi\|\le\epsilon$ for a
    unit vector $\psi$, then
    \begin{align}
        \|(B_j-\langle\psi,B_j\psi\rangle)\psi\|
           &\le4\ell\epsilon,&
        \|B_j-\langle\psi,B_j\psi\rangle\|
           &\le8\ell J.
        \label{eq:hastings05_rounded_pairs}
    \end{align}
\end{lemma}
\begin{proof}\leanok
    A radius-$\ell$ ball assigned to a nonclosing edge meets only its
    two adjacent blocks. The closing assignment collects the two end
    intervals. Counting their centres gives the fibre bound $4\ell$;
    summing over fibres proves both estimates. For a short chain $N=1$,
    the same assignment puts every centre on the closing edge.
\end{proof}

\begin{definition}[The constructed spectral hierarchy]
    \label{def:hastings05_constructed_hierarchy}
    \lean{TNLean.MPS.AreaLaw.LocalSolution.hierarchyEdge,
        TNLean.MPS.AreaLaw.LocalSolution.extendFiniteEdges,
        TNLean.MPS.AreaLaw.LocalSolution.finiteHierarchyDepth,
        TNLean.MPS.AreaLaw.LocalSolution.hierarchyProjection,
        TNLean.MPS.AreaLaw.LocalSolution.hierarchyClosingProjection,
        TNLean.MPS.AreaLaw.LocalSolution.hierarchyWindowRadius,
        TNLean.MPS.AreaLaw.LocalSolution.periodicHierarchyProjection,
        TNLean.MPS.AreaLaw.LocalSolution.roundedPeriodicProjection,
        TNLean.MPS.AreaLaw.LocalSolution.IsRetainedHierarchy}
    \leanok
    Extend the finite edge family by zero and set
    $A_k^m=B_{(2k+1)2^m}$. Start with $P_k^0=I$. Given the two child
    projections, let $Q=P_{2k}^mP_{2k+1}^m$ and define
    \begin{align}
        P_k^{m+1}
          &=Q\,\mathbf1_{[c_k^m-x_m,c_k^m+x_m]}(QA_k^mQ),&
        c_k^m&=\langle\psi,A_k^m\psi\rangle.
        \label{eq:hastings05_constructed_window}
    \end{align}
    Put $h=\max(2,\lceil\log_2N\rceil)$. After the $h$ binary levels,
    apply one further compressed window to $B_0$ inside
    $\operatorname{ran}P_0^h$. Denote the resulting projection by $P$.
    For $K>0$, $r\ge0$, the radii are
    $x_m=K(e^{-r}+2\beta_m)^{1/h}$, where
    $\beta_m$ is the majorant of
    Definition~\ref{def:hastings05_hierarchy_majorant} with
    $e=e^{-r}$ and $s=1-1/h$.
\end{definition}

\begin{lemma}[The chosen windows and their descendants]
    \label{lem:hastings05_constructed_windows}
    \lean{TNLean.MPS.AreaLaw.LocalSolution.norm_sub_hierarchyProjection_le,
        TNLean.MPS.AreaLaw.LocalSolution.exists_unit_mem_hierarchyClosingProjection,
        TNLean.MPS.AreaLaw.LocalSolution.hierarchyDeficiency_window_succ_le,
        TNLean.MPS.AreaLaw.LocalSolution.hierarchyWindowRadius_le,
        TNLean.MPS.AreaLaw.LocalSolution.hierarchyProjection_children,
        TNLean.MPS.AreaLaw.LocalSolution.isRetainedHierarchy_of_hierarchyProjection,
        TNLean.MPS.AreaLaw.LocalSolution.isRetainedHierarchy_of_hierarchyClosingProjection}
    \leanok
    \uses{def:hastings05_constructed_hierarchy,
        thm:hastings05_hierarchy_closed_form,thm:hastings05_supersite_hierarchy}
    Suppose every edge has residual at most $Ke^{-r}$ and centred norm
    at most $K$. If $h\ge2$, then the deficiency after the closing
    window is at most $(3^{h+1}-1)e^{-r/8}$. Whenever this is strictly
    less than one, $P$ retains a unit vector within twice that distance
    of $\psi$. Every vector retained by $P$ is retained at each node
    of the underlying binary hierarchy. For $m\le h$,
    $x_m\le9K e^{-r/(4h)}$.
\end{lemma}
\begin{proof}\leanok
    Disjoint child supports imply commutation. Apply the deficiency
    recurrence to the recursively defined compressed windows, including
    the closing window, and normalize the nonzero image of $\psi$.
    A parent window has range in the product of its child ranges;
    induction proves membership at every descendant.
\end{proof}

\begin{lemma}[The gap of the original Hamiltonian]
    \label{lem:hastings05_physical_energy_comparison}
    \lean{TNLean.MPS.AreaLaw.LocalSolution.abs_re_inner_sub_le_norm_sub,
        TNLean.MPS.AreaLaw.LocalSolution.card_le_two_mul_re_trace_of_approximation,
        TNLean.MPS.AreaLaw.LocalSolution.card_le_two_mul_re_trace_of_periodicHierarchyProjection,
        TNLean.MPS.AreaLaw.LocalSolution.card_le_two_mul_re_trace_of_roundedPeriodicProjection}
    \leanok
    Suppose $H\ge E_0I+\Delta(I-P_0)$, where $P_0$ is an orthogonal
    projection, and $\psi$ is a unit ground vector. Let $T$ satisfy
    $\|H-T\|\le\delta$. If every unit vector in a retained subspace
    has $T$-energy at most $\langle\psi,T\psi\rangle+\nu$ and
    $2\delta+\nu\le\Delta/2$, that subspace has dimension at most
    $2\operatorname{tr}P_0$.
\end{lemma}
\begin{proof}\leanok
    Each unit vector in the subspace has $H$-energy at most
    $E_0+2\delta+\nu$, hence ground-projection expectation at least
    $1/2$. Sum this inequality over an orthonormal family and bound
    the sum of its projection expectations by $\operatorname{tr}P_0$.
\end{proof}

\begin{theorem}[The physical periodic hierarchy under local residuals]
    \label{thm:hastings05_physical_hierarchy}
    \lean{TNLean.MPS.AreaLaw.LocalSolution.periodicPhysicalTruncatedTerm,
        TNLean.MPS.AreaLaw.LocalSolution.exists_periodicPhysicalTruncatedTerm_localization,
        TNLean.MPS.AreaLaw.LocalSolution.physicalPeriodicProjectionRadius,
        TNLean.MPS.AreaLaw.LocalSolution.physicalPeriodicProjectionRadius_spec,
        TNLean.MPS.AreaLaw.LocalSolution.physicalPeriodicProjectionRadius_lt,
        TNLean.MPS.AreaLaw.LocalSolution.exists_periodicPhysicalHierarchy}
    \leanok
    \uses{thm:hastings05_metric_filter,lem:hastings05_zero_range,
        lem:hastings05_rounded_pairs,lem:hastings05_constructed_windows,
        lem:hastings05_physical_energy_comparison}
    Fix $R\ge0$, $J>0$ and $\Delta>0$. There are $c_1,\eta>0$ and
    $C\ge0$, chosen before the chain length and interactions, as follows.
    Let $H=\sum_i g_i$ on a nonempty periodic chain of $M$ sites, where
    $g_i=g_i^\dagger$ is supported within cyclic distance $R$ of $i$
    and $\|g_i\|\le J$. Let $\psi$ be a unit ground vector of energy
    $E_0$, and suppose $H\ge E_0I+\Delta(I-P_0)$ for an orthogonal
    projection $P_0$. Put $J_b=\max(J,JC)$, $L=\lceil\log_2M\rceil+2$,
    and choose
    \begin{align}
        \ell&=\max\left(\lceil R\rceil,
          \left\lceil\frac{512L^2}{\eta}
          \left(1+\max\left(0,\log\frac{8J_b}{\eta\Delta}\right)\right)
          \right\rceil\right),&r&=\eta\ell.
        \label{eq:hastings05_physical_radius}
    \end{align}
    Construct the actual truncated filters $a_i=M_i$ by
    (\ref{eq:hastings05_metric_filter}), and the rounded pair family
    and projection $P$ by Lemma~\ref{lem:hastings05_rounded_pairs} and
    Definition~\ref{def:hastings05_constructed_hierarchy}, with
    $K=8\ell J_b$. Assume the local residual condition
    \begin{align}
        \|(a_i-\langle\psi,a_i\psi\rangle)\psi\|
           &\le J_b e^{-r}.
        \label{eq:hastings05_physical_low}
    \end{align}
    Then $P$ retains a unit vector $\phi$ satisfying
    $\|\phi-\psi\|\le2(3^{h+1}-1)e^{-r/8}$. This vector belongs to
    the retained space at every node of the constructed hierarchy.
    Every orthonormal family retained by $P$ has at most
    $2\operatorname{tr}P_0$ members.
\end{theorem}
\begin{proof}\leanok
    For positive $R$, apply Theorem~\ref{thm:hastings05_metric_filter}
    with the uniform cyclic-ball bound $2\lceil R\rceil+1$. For $R=0$,
    use Lemma~\ref{lem:hastings05_zero_range}. The rounded fibre count
    gives all pair residual and norm estimates. The radius in
    (\ref{eq:hastings05_physical_radius}) makes the sum of the binary
    and closing window errors at most $\Delta/4$ and the truncation
    error at most $\Delta/8$. Apply
    Lemmas~\ref{lem:hastings05_constructed_windows} and
    \ref{lem:hastings05_physical_energy_comparison}.
\end{proof}

The local residual assumption
(\ref{eq:hastings05_physical_low}) is the hypothesis used in
\cite[eq.~(low), lines 493--507 and 633--639]{Hastings2006Solving}.
In a degenerate ground space, small leakage outside that space does not
by itself give this scalar residual for an arbitrarily chosen ground
vector. A simultaneous choice of ground vectors satisfying the required
condition is a separate question
\cite{gap:hastings05_two_dim_ground_basis}. The nesting conclusion
describes retained subspaces; no bound for intermediate tensor dimensions
is asserted here.

\begin{corollary}[The hierarchy for a specified unique ground state]
    \label{cor:hastings05_unique_physical_hierarchy}
    \lean{TNLean.MPS.AreaLaw.LocalSolution.exists_periodicPhysicalHierarchy_of_isGappedGroundState}
    \leanok
    \uses{thm:hastings05_metric_filter,lem:hastings05_zero_range,
        lem:hastings05_rounded_pairs,lem:hastings05_constructed_windows,
        lem:hastings05_physical_energy_comparison}
    Fix $R\ge0$, $J>0$ and $\Delta>0$. There are $c_1,\eta>0$ and
    $C\ge0$, chosen before the physical chain and all ground-state data,
    with the following property. Let $H=\sum_i g_i$ on a nonempty
    periodic chain, with Hermitian terms of norm at most $J$ supported
    within cyclic distance $R$ of their centres. Let $\Omega$ be a
    specified unit vector such that
    $H\Omega=E_0\Omega$ and
    $H\ge E_0I+\Delta(I-|\Omega\rangle\langle\Omega|)$.
    Put $J_b=2\max(J,JC)$ and choose the radius by
    (\ref{eq:hastings05_physical_radius}). The actual filters, rounded
    pair family and spectral windows for this ground vector retain a
    unit vector $\phi$ satisfying
    $\|\phi-\Omega\|\le2(3^{h+1}-1)e^{-\eta\ell/8}$.
    This vector belongs to every descendant retained space, and the
    final retained space has dimension at most two. No local residual
    estimate is assumed.
\end{corollary}
\begin{proof}\leanok
    Start with the localization constants $c_1,\eta_{\rm loc},C$ and
    replace the rate by $\eta=\min(\eta_{\rm loc},c_1\Delta/2)$.
    The unique-ground Gaussian estimate gives one spectral leakage
    term and one localization error. Consequently
    \begin{align}
        \|(a_i-\langle\Omega,a_i\Omega\rangle)\Omega\|
        &\le J e^{-c_1\ell\Delta/2}+JC e^{-\eta_{\rm loc}\ell}
         \le (J+JC)e^{-\eta\ell}
         \le J_b e^{-\eta\ell}.
        \label{eq:hastings05_unique_physical_low}
    \end{align}
    The same constants bound the localization error and the term norm.
    Apply the constructed-window estimate at the chosen radius. The
    rank-one ground projection has trace one, so
    Lemma~\ref{lem:hastings05_physical_energy_comparison} bounds the
    retained dimension by two.
\end{proof}

The filters and the assignment of their centres to pairs are determined
by the local interactions and the chosen parameters. The spectral-window
centres are expectations in the specified vector $\Omega$. This is a
representation argument for that ground state, rather than an algorithm
that constructs the windows without ground-state data.
